\subsection{The complete next ordinary class groups from the actual norm lattice}

\begin{statement}{E.106}{Proved}{The ninth-symbol transfer, vanishing transvection, and all third-level class maps}
Retain the literal roots, prime labels and unit basis of E.76,
E.86, E.91 and E.101:
\[
 K=\mathbb Q(\zeta),\quad \xi^3=\zeta,\quad
 L_d=K(\beta_d),\quad B_d=L_d(\xi),\qquad
 \beta_{d+1}^3=\beta_d,\quad\beta_1^3=17\zeta.
 \tag{E106.1}
\]
Write $E_F=\mathcal O_F^\times$ and put a hat on a unit group
for its completion for multiplication by $3^m$. The first nine
E.101 generators of $E_{B_1}$ are $u_1=\xi,u_2,\ldots,u_9$;
their only three-primary relation is $u_1^9=1$.
The complete ordinary norm image from $B_2$ has the specified coordinates
\[
 \begin{aligned}
 N_{B_2/B_1}\widehat E_{B_2}
 &=\left\{\prod_{i=1}^9u_i^{x_i}:x=P y\right\},\\
 Py&=(3y_1+y_2+y_4,\ y_2,\ y_3,\ y_4,
             3y_5+2y_6,\ y_6,\ y_7,\ 3y_8,\ 3y_9)^t,\\
 y&\in\mathbb Z_3^9/3\mathbb Z_3e_1 .
 \end{aligned}
 \tag{E106.2}
\]
For $v\in\widehat E_{B_2}$, take the actual $y$-coordinates
of $N_{B_2/B_1}v$ in (E106.2). Its six arithmetic cubic
local norm characters for $B_3/B_2$, ordered by the three
primes over three and then the three over seventeen, are
\[
 \chi(v)=D y\pmod3,\qquad
 D=\begin{pmatrix}
 2&1&2&2&0&2&2&2&1\\
 2&1&2&0&0&2&0&0&0\\
 2&1&2&1&0&2&1&1&2\\
 1&2&1&0&0&0&2&2&2\\
 1&2&1&1&0&0&0&2&2\\
 1&2&1&2&0&0&1&2&2
 \end{pmatrix}.
 \tag{E106.3}
\]
The arithmetic generator sends $\beta_3$ to $\zeta\beta_3$.
The image of this actual character map has dimension four.
Its kernel is the completed ordinary-unit subgroup of field
norms from $B_3$. The certificate gives the full finite row
and column transformations, their inverses, a four-column
image basis, and its coordinate inverse.

The entire ordinary three-class groups are
\[
 \begin{gathered}
 \operatorname{Cl}(B_3)\{3\}
       \xrightarrow[\mathrm{Art}]{\ \sim\ }
       A_9=C_9^2\oplus C_3,\qquad
 W=(g_2-1)\operatorname{Cl}(B_3)\{3\}=0,\\
 \operatorname{Cl}(L_3)\{3\}=C_9a_3\oplus C_3b_3,\qquad
 [\mathfrak q_{3,3}]=3a_3,\qquad
 \operatorname{ord}[\mathfrak q_{3,3}]=3,\\
 N_{L_3/K}\widehat{\mathcal O_{L_3}[1/3]^\times}_3
                   =\lambda^{3\mathbb Z_3},\qquad
 a_3^{\rm E69}=1 .
 \end{gathered}
 \tag{E106.4}
\]
Here $a_3,b_3$ are the literal coherent ideal classes of E86.4;
$a_3^{\rm E69}$ distinguishes the projection exponent from
that ideal class. The first arrow is the existing restriction
to the actual ninth-genus field. That field is therefore the
full Hilbert three-class field of $B_3$.

Use $(y,z,t)$ for the $A_9$ coordinates, $(s,u,v)$ for the
$C_9\oplus C_3^2$ coordinates of $\operatorname{Cl}(B_2)\{3\}$,
and $(a_d,b_d)$ for the specified $L_d$ coordinates. All
adjacent and horizontal ordinary class maps at this level are
\[
 \begin{array}{c|c|c}
 \text{extension}&\text{ideal norm}&\text{ideal extension}\\ \hline
 B_3/B_2&
 \begin{pmatrix}2&3&6\\2&0&2\\0&1&0\end{pmatrix}&
 \begin{pmatrix}6&0&0\\0&0&3\\0&0&0\end{pmatrix}\\[3pt]
 L_3/L_2&I_2&3I_2\\[3pt]
 B_3/L_3&
 \begin{pmatrix}3&6&6\\1&0&0\end{pmatrix}&
 \begin{pmatrix}3&0\\3&3\\2&0\end{pmatrix}.
 \end{array}
 \tag{E106.5}
\]
Every row is reduced modulo its target cyclic order.
The vertical $B$ norm kernel is $\langle(0,3,0)\rangle=C_3$.
The horizontal extension kernel is precisely
$\langle3a_3\rangle=\langle[\mathfrak q_{3,3}]\rangle$.
The actual ordinary unit norm is surjective:
\[
                  N_{B_3/L_3}E_{B_3}=E_{L_3}.
 \tag{E106.6}
\]

The six actual S-prime classes are now completely determined:
\[
 C_S=
 \begin{array}{c|rrrrrr}
 &Q_1&Q_2&Q_3&T_1&T_2&T_3\\ \hline
 C_9&6&6&6&6&3&0\\
 C_9&0&6&3&5&8&5\\
 C_3&2&2&2&1&1&1
 \end{array}.
 \tag{E106.7}
\]
They generate the order-$81$ subgroup
$R_3=3C_9\oplus C_9\oplus C_3$. Its actual quotient map is
$(y,z,t)\mapsto y\bmod3$, in the coherent prime-$73$ S-basis.
All three $Q_i$ have order three, and their product is principal.
The full principal S-divisor lattice has basis and index
\[
 \mathcal L_3=H_S\mathbb Z^6,\qquad
 H_S=\begin{pmatrix}
 3&1&-3&2&2&0\\
 0&1&0&0&0&0\\
 0&-2&3&-2&-2&0\\
 0&0&0&3&0&1\\
 0&0&0&0&3&1\\
 0&0&0&0&0&1
 \end{pmatrix},\qquad
 [\mathbb Z^6:\mathcal L_3]=81 .
 \tag{E106.8}
\]
For a divisor vector $(q_1,q_2,q_3,t_1,t_2,t_3)$ the inverse
coordinates are
\[
 \left(
 {q_1+q_2+q_3\over3},\ q_2,\
 {q_3+2q_2+2B+2C\over3},\ B,\ C,\ t_3
 \right),\quad
 B={t_1-t_3\over3},\quad C={t_2-t_3\over3}.
 \tag{E106.9}
\]
Thus an integral S-divisor is principal exactly when all
these coordinates are integral. The same criterion holds
after completion, with integrality in $\mathbb Z_3$.

There is an exact complete ordinary-unit norm test from
$B_3$ to $B_2$. First apply (E106.3). On its zero kernel,
for an ordinary unit $v$ choose a field norm preimage
$w\in B_3^\times$ and solve
$\operatorname{div}(w)=(\sigma_2-1)\mathfrak D$.
Write $\mathrm{Art}[\mathfrak D]_3=(y,z,t)$ in (E106.4).
Then
\[
 \begin{gathered}
 v\in N_{B_3/B_2}E_{B_3}
   \quad\Longleftrightarrow\quad
   \chi(v)=0\text{ and }y(\mathrm{Art}[\mathfrak D]_3)=0\pmod3,\\
 E_{B_2}/N_{B_3/B_2}E_{B_3}\simeq C_3^5,\qquad
 \ker(N:\widehat E_{B_3}\to\widehat E_{B_2})
                     \simeq C_3\oplus\mathbb Z_3^{54}.
 \end{gathered}
 \tag{E106.10}
\]
The second test is independent of both choices, and both
directions are proved below. On completion it is the
continuous extension of this finite quotient test. Each
nonempty completed norm fiber is the coset of the actual
norm kernel through any of its preimages. The corresponding
$L_3/L_2$ ordinary norm quotient is $C_3^3$, with norm kernel
$C_3\oplus\mathbb Z_3^{18}$.

These conclusions evaluate the third finite projection and
the E.91 transvection. The eventual exponent $h$ of E.69
and the later ordinary prime orders are not assigned a
value by this finite computation. The manuscript anchor is
LM~\S51.16, folios 733v--736v, scans 172--178, through the actual
positional norm maps of E.46, E.69, E.72 and E.88.
\end{statement}

\begin{proof}
\emph{The local norm square and its precise exponent.}
At each of the six S-primes of $B_1$ there is one prime in
each vertical upper field, with residue degree one.
The local extension $B_3/B_1$ is cyclic of degree nine:
$\mu_9\subset B_1$, $\beta_3^9=\beta_1$, and its generator
$g$ sends $\beta_3$ to $\xi\beta_3$. Its subgroup over
$B_2$ is generated by $g^3=\sigma_2$.
The specified local Artin maps give the commutative square
\[
 \begin{array}{ccc}
 B_{2,w}^{\times}&\xrightarrow{\ \mathrm{Art}\ }&
             \operatorname{Gal}(B_{3,W}/B_{2,w})\\
 \big\downarrow\scriptstyle N&&\big\downarrow\scriptstyle\mathrm{inclusion}\\
 B_{1,v}^{\times}&\xrightarrow{\ \mathrm{Art}\ }&
             \operatorname{Gal}(B_{3,W}/B_{1,v}).
 \end{array}
 \tag{E106.11}
\]
The right arrow takes $\sigma_2^j$ to $g^{3j}$. Hence
\[
       \chi_{9,v}(Nu)=3\chi_{3,w}(u)\quad\text{in }\mathbb Z/9.
 \tag{E106.12}
\]
This factor three is the literal inclusion
$\mu_3\hookrightarrow\mu_9$, $\zeta^j=\xi^{3j}$.
With the norm variable first, the equivalent Hilbert-symbol
identity is
$(u,\beta_2)_{B_{2,w},3}
 =(u,\beta_1)_{B_{2,w},9}
 =(Nu,\beta_1)_{B_{1,v},9}$.
The Artin square fixes the orientation regardless of
which Hilbert-symbol argument convention is used.
It follows from the restriction/corestriction compatibility
of the proved Tate cup-product isomorphism
\cite[III,Theorem3.1 and its following proof; III,\S4]{E_MilneCFT}.
This uses that proved compatibility; III,Lemma3.2 itself
prints that its proof is to be added.

The E.76 ninth-character rows on the literal nine units are
\[
 M_9=\begin{pmatrix}
 1&5&3&2&6&0&3&4&2\\
 1&5&3&8&3&6&0&3&6\\
 1&5&3&5&0&3&6&8&7\\
 5&7&6&4&0&0&3&1&4\\
 8&4&6&7&3&3&0&1&1\\
 1&5&3&5&3&3&3&5&8
 \end{pmatrix}\pmod9.
 \tag{E106.13}
\]
The exact arithmetic normalization of E94.13 multiplies
these rows by $(5,5,5,8,5,4)$ modulo nine. Their normalized
sum is zero on every global unit. The certificate checks
the equality of all nine literal unit polynomials in E.76
and E.101 before composing these maps.

E.101 proves that $P$ is the complete ordinary norm-image
lattice. Each of its nine generators is an actual ordinary
norm. Every ordinary norm has these coordinates, with the
sole relation $3e_1$. Equation (E106.12) consequently proves
both containments in the character image, and gives
\[
 D=\left({1\over3}\operatorname{diag}(5,5,5,8,5,4)M_9P\right)\bmod3.
 \tag{E106.14}
\]
Every entry before division is divisible by three. Changing
a ninth-character representative changes the divided entry
by a multiple of three. The torsion relation is killed.
The complete finite row and column transformations reduce
$D$ to a rank-four identity block; their inverses are
retained. Enumeration of all $3^9$ source residues gives
$81$ image vectors, each with $243$ preimages.

Thus the local field-norm index on ordinary units is
exactly $81$. Units are norms at unramified finite primes,
and all infinite primes are complex. The cyclic Hasse
norm theorem gives the asserted global kernel
\cite[VIII,Theorem3.1 and proof]{E_MilneCFT}.
The same finite test holds on completion: the unit group
is dense and the finite character map is continuous.
Equivalently, the completed norm image in the full
multiplicative group is the kernel of the completed
cyclic norm-obstruction map of E.64.

\emph{The entire class group and the prime product.}
There are precisely six ramified finite primes in
$B_3/B_2$, all of index three. E.91 proves that the vertical
generator acts identically on the entire ordinary three-class
group. The ambiguous class number formula therefore computes
the full order:
\[
 |\operatorname{Cl}(B_3)\{3\}|
  =|\operatorname{Cl}(B_2)\{3\}|\,{3^{6-1}\over3^4}
  =81\cdot3=243 .
 \tag{E106.15}
\]
The denominator is the ordinary-unit \emph{field}-norm
index, not the ordinary-unit norm index
\cite[Theorem1(1), with its ideal/Herbrand derivation]{E_LemmermeyerAmbiguous}.
The actual Artin map of E.91 is onto $A_9$, also of order
$243$, with kernel $W$. Thus it is an isomorphism and $W=0$.
Its existing prime columns give the full classes (E106.7).

Their subgroup has order $81$. The determinant of $H_S$
is $81$, and its columns map to zero under $C_S$ with
target moduli $(9,9,3)$. Equality of the indices proves
that it is the full kernel. Solving its six equations
gives (E106.9). Each S-prime class is three-primary:
a suitable power extends the corresponding B1 S-prime,
whose class has order dividing three. Hence this kernel
is actual principality for S-divisors. In particular
$Q_1Q_2Q_3$ and every $Q_i^3$ are principal.

The horizontal cubic $B_3/L_3$ is the actual S-split
unramified extension of E.86. Its prime norms send each
$Q_i$ to $\mathfrak q_{3,3}$. Thus $\mathfrak q_{3,3}^3$
is principal. Its norm to $L_2$ has order three, proving
that its order is exactly three. E.91's group-type theorem
now gives $\operatorname{Cl}(L_3)\{3\}=C_9\oplus C_3$.
Norm to $L_2$ is onto between groups of order $27$, so is
an isomorphism; it sends the literal coherent $a_3,b_3$
to $a_2,b_2$. Therefore $[\mathfrak q_{3,3}]=3a_3$.
The exact valuation/norm comparison in E.72 identifies
this order with the E.69 finite projection exponent.

\emph{All ordinary norm and extension matrices.}
The vertical $B$ norm uses the actual ray restriction of
E.76. Write $U_3,U_9$ for its integral Smith left maps,
$R$ for its ray restriction matrix, $\iota$ for E91.7's
simple inclusion, and $p$ for E86.7's simple projection.
The literal composition is
\[
                         N=pU_3R U_9^{-1}\iota .
 \tag{E106.16}
\]
Its cyclic reduction is the first matrix of (E106.5).
It is onto and its kernel is exactly
$\{(0,0,0),(0,3,0),(0,6,0)\}$. Since the vertical class
action is identity, ideal extension obeys $jN=3$.
Three times any norm preimage is independent of its
choice, yielding the displayed matrix $j$. The certificate
checks both identities $jN=3$, $Nj=3$ on every class and
checks all six specified prime norms.

The $L$ norm is the isomorphism already proved. Its
equivariance, $N\sigma=N$, makes the vertical class action
identity, hence its extension is $3I_2$. Commuting the
actual field norms in $(B_3,B_2,L_3,L_2)$, with E86.9,
gives the horizontal norm matrix.

For the horizontal extension the actual primes remove
the possible lift ambiguity. The coherent prime above
$73$ has class $(1,8,2)$; its norm is $J_3$. The
extension of $b_3=[J_3]_3$ is consequently its $h$-trace:
\[
 H=\begin{pmatrix}4&3&0\\8&7&6\\1&0&1\end{pmatrix},
 \qquad (1+H+H^2)(1,8,2)^t=(0,3,0)^t\quad\text{in }A_9.
 \tag{E106.17}
\]
The prime $P_3=(7,\beta_3+1)$ is inert in $B_3/L_3$:
its residue field does not contain $\mu_9$. Its extended
prime has residue field $\mathbb F_{343}$. Its norm to
B1 is $(7,\beta_1+1)$, with residue degree one in that
vertical comparison. This B1 prime splits in $B_3/B_1$
because $\beta_1=-1$ is a ninth power in its residue field.
The fresh certified B1 ray calculation gives
\[
 \mathrm{Art}[(7,\beta_1+1)]=(6,0,1),\qquad
 \mathrm{Art}[(73,\xi-55,\beta_1-22)]=(1,8,2).
 \tag{E106.18}
\]
Both complete ideal HNFs and ray/Smith coordinates are
retained. Arithmetic Frobenius restriction is the actual
norm comparison. Since $a_3=[J_3/P_3]_3$, subtracting
the first column in (E106.18) from (E106.17) yields
$j(a_3)=(3,3,2)$, $j(b_3)=(0,3,0)$.
The identities $jN=1+H+H^2$, $Nj=3$ hold on every class,
and $\ker j=\langle3a_3\rangle$.

For this unramified cubic extension the invariant
fractional ideals are exactly extended base ideals.
Hilbert--90 and the principal-ideal sequence identify the
capitulation kernel with $H^1(C_3,E_{B_3})$.
The unit Herbrand quotient is $1/3$, since all infinite
primes are complex. Therefore
\[
                  |\ker j|=3[E_{L_3}:N_{B_3/L_3}E_{B_3}].
 \tag{E106.19}
\]
Its left side is three. This proves (E106.6) for the
actual unit norm. These are the ideal and Herbrand maps
used in E.86 and in the cited ambiguous-class derivation.

\emph{The remaining ordinary-unit defect and completion.}
For $B_3/B_2$, every invariant ideal is a product of an
extended base ideal and powers of the six ramified primes.
Thus its strongly ambiguous three-class subgroup is
$j\operatorname{Cl}(B_2)\{3\}+R_3=R_3$, while the full
ambiguous group is $A_9$. For a unit that is a field norm,
the ideal Hilbert--90 equation in the statement exists,
orbit by orbit on primes. Changing its solution adds an
invariant ideal. Changing the norm preimage changes it
by a principal divisor followed by that invariant ambiguity.
Its exact choice-independent obstruction is
\[
                 A_9/R_3\xrightarrow{\ \sim\ }C_3,\qquad
                 (y,z,t)\longmapsto y\bmod3 .
 \tag{E106.20}
\]
Vanishing lets the divisor be adjusted to a principal one;
the resulting Hilbert--90 correction makes the norm preimage
an ordinary unit. An ordinary unit preimage conversely has
zero divisor. Every ambiguous ideal class yields a field
norm by this construction, proving surjectivity of the
extra obstruction. This proves both directions in
(E106.10), with the actual ideal-class map
\cite[Proposition1]{E_LemmermeyerAmbiguous}.

The extra index three times the local index $81$ gives
ordinary norm index $243$. Inclusion of a base unit has
norm its cube, so this quotient is killed by three and is
$C_3^5$. For $L_3/L_2$, the now trivial class action and
its two ramified primes give field-unit norm index three
by the ambiguous formula. Its strongly ambiguous subgroup
is $3D_3=\langle\mathfrak q_{3,3}\rangle$, of order three;
the other ramified prime is principal. Its additional
defect has order nine. The ordinary norm quotient is
therefore $C_3^3$.

Each norm image contains the cubes and has finite cokernel.
Completion of these finitely generated unit groups preserves
every displayed image and finite test. The kernel free ranks
are $80-26=54$ for B and $26-8=18$ for L. Their torsion
kernels are the kernels of cubing on $\mu_9$ or $\mu_3$,
respectively, namely $\mu_3$. This proves the fiber groups.
Unique-prime, residue-degree-one valuation maps also show
that an S-unit of norm one has zero S-valuations.

The computation uses the complete E.101 norm image.
All finite transformations, the $19\,683$ character inputs,
all $243$ upper and $81$ lower class coordinates, and the
two actual prime Frobenii are retained. It uses no BNF
computation in $B_2,B_3,L_2$ or $L_3$.
\end{proof}
